\documentclass[11pt]{article}

\usepackage[a4paper,margin=1in]{geometry}
\usepackage{amsmath,amssymb,amsthm,mathtools}
\usepackage{booktabs}
\usepackage{enumitem}
\usepackage[hidelinks]{hyperref}
\usepackage{xcolor}

\title{\vspace{-1.5em}Corrector Timing as Statistical Intervention Diagnosis\\
\large A supervisor-facing theory memo}
\author{Matteo Omizzolo}
\date{May 2026}

\newtheorem{definition}{Definition}
\newtheorem{proposition}{Proposition}
\newtheorem{lemma}{Lemma}
\newtheorem{theorem}{Theorem}
\newtheorem{corollary}{Corollary}
\newtheorem{remark}{Remark}

\newcommand{\calX}{\mathcal X}
\newcommand{\calT}{\mathcal T}
\newcommand{\calS}{\mathcal S}
\newcommand{\calF}{\mathcal F}
\newcommand{\E}{\mathbb E}
\newcommand{\R}{\mathbb R}
\newcommand{\1}{\mathbf 1}
\newcommand{\TV}{\mathrm{TV}}
\newcommand{\KL}{\mathrm{KL}}
\newcommand{\osc}{\mathrm{osc}}
\newcommand{\Var}{\mathrm{Var}}
\newcommand{\Cov}{\mathrm{Cov}}

\begin{document}
\maketitle

\begin{abstract}
This memo locks the thesis direction around a statistical, rather than
algorithmic, theory of corrector timing.  The central object is the value of a
corrector call as an intervention in an inhomogeneous generative Markov chain.
The proposed contribution is a finite-budget diagnostic framework: define
corrector-timing estimands, decompose schedule gains into marginal and
interaction effects, state safe regret implications when low-order structure is
present, and connect the framework to informed correctors, locally balanced
MCMC proposals, and non-uniform scan Gibbs.  The mathematical claims here are
deliberately modest.  We do not claim that ProSeCo is locally balanced, that the
schedule objective is submodular, or that a search procedure is optimal.  The
framework is model- and corrector-agnostic; ProSeCo-OWT is one empirical
instantiation used to diagnose a non-marginal, redundant, locally exploitable
timing regime.
\end{abstract}

\section{Direction lock}

The thesis should not be framed as ``we designed a better schedule optimizer.''
That is too close to a model-specific engineering problem.  The more stable
statistical question is:

\begin{quote}
For a fixed masked-diffusion predictor, fixed informed corrector, quality
functional \(F\), and finite corrector-call budget \(B\), can we diagnose when
corrector timing is marginal, state-dependent, redundant, or interaction-driven?
\end{quote}

Algorithms then become probes.  A separable ranker probes whether marginal
effects are sufficient.  Pairwise diagnostics probe whether interactions are
structured and composable.  Local or true-\(G\) search probes whether useful
structure remains after low-order summaries fail.  Bayesian optimization is a
possible downstream method only if the diagnostics reveal a smooth kernel
structure; it is not a premise of the thesis.

\subsection{The gap}

Existing work answers adjacent but different questions.

\begin{itemize}[leftmargin=1.5em]
\item Informed-corrector work for discrete diffusion designs better local
corrector transitions.  Zhao, Shi, Chen, Druckmann, Mackey, and Linderman
propose predictor-corrector schemes where the corrector uses model information,
including hollow-transformer modifications, to counter approximation error in
discrete diffusion sampling \cite{zhao2024informed}.
\item Locally balanced proposals answer how local MCMC proposals in discrete
spaces should use local target information.  Zanella characterizes locally
balanced proposals and proves Peskun-type optimality in high-dimensional
regimes \cite{zanella2020informed}.
\item Random-scan, non-uniform, and adaptive Gibbs work studies how update
schemes or selection probabilities affect MCMC convergence
\cite{roberts1997updating,chimisov2018adapting}.
\item Lavenant and Zanella study error bounds and optimal predictor schedule
sizes for masked diffusions with factorized approximations \cite{lavenant2025}.
\end{itemize}

The thesis gap is the timing layer:

\begin{quote}
Given a fixed predictor and a fixed informed corrector kernel, when should a
limited number of corrector calls be deployed along the generation trajectory?
\end{quote}

This is not a new corrector-kernel theorem and not a new unmasking-schedule
theorem.  It is a finite-budget intervention problem over an inhomogeneous
sequence of local correction opportunities.

\subsection{Candidate theory options}

There are three plausible mathematical directions.  They are not equally safe.

\begin{center}
\begin{tabular}{p{0.22\linewidth}p{0.34\linewidth}p{0.32\linewidth}}
\toprule
Option & What it would prove & Assessment \\
\midrule
Kernel theorem & ProSeCo or a Zhao-style corrector is locally balanced,
contractive, or Peskun-optimal & Too strong for this thesis unless the kernel
is redesigned and analysed directly \\
Optimizer theorem & A particular search or greedy scheduler is optimal or
near-optimal & Not aligned with the diagnostic/statistical direction; risks
becoming model-specific optimization \\
Diagnostic theory & Corrector timing has well-defined intervention effects,
exact interaction decompositions, and safe low-order implications & Most
promising: mathematically clean, corrector-agnostic, and directly connected to
the experiments \\
\bottomrule
\end{tabular}
\end{center}

The rest of this memo develops the third option.  The stronger kernel and
optimizer theories appear only as intuition or future work.

\section{Mathematical setup}

Assume a finite or countable state space \(\calX\).  A masked-diffusion sampler
with fixed predictor is represented as an inhomogeneous Markov chain
\[
  X_0 \sim \mu_0,\qquad
  X_t \sim P_t(X_{t-1},\cdot), \qquad t=1,\ldots,T,
\]
where \(P_t\) is the predictor transition at denoising step \(t\).  A fixed
corrector at time \(t\) is a Markov kernel \(K_t\) on \(\calX\).  A schedule is
a subset \(S\subseteq [T]:=\{1,\ldots,T\}\).  Under schedule \(S\), the law of
the terminal sample is
\[
  \mu_T^S
  =
  \mu_0 P_1 K_1^{\1\{1\in S\}}
        P_2 K_2^{\1\{2\in S\}}
        \cdots
        P_T K_T^{\1\{T\in S\}},
\]
where \(K_t^0\) is the identity kernel and \(K_t^1=K_t\).  Let
\(F:\calX\to\R\) be a bounded terminal quality functional.  Define
\[
  J(S) := \E_{\mu_T^S}[F(X_T)],
  \qquad
  G(S) := J(S)-J(\varnothing).
\]
The fixed-budget timing problem considers schedules with \(|S|=B\).

\begin{remark}[What is abstracted away]
This notation does not require \(K_t\) to be Gibbs, Metropolis-within-Gibbs, or
locally balanced.  It only requires that the corrector call can be represented
as a kernel inserted into the trajectory.  This is why the framework is
corrector-agnostic.
\end{remark}

\begin{definition}[Timing estimands]
For \(t\in[T]\), define the marginal intervention effect
\[
  \Delta_t := G(\{t\}).
\]
For \(s\ne t\), define the pair interaction
\[
  \xi_{s,t}
  :=
  G(\{s,t\}) - G(\{s\}) - G(\{t\}).
\]
For \(A\subseteq B\subseteq [T]\) and \(x\notin B\), define the diminishing
returns gap
\[
  \mathrm{DR}(A,B,x)
  :=
  \bigl[G(A\cup\{x\})-G(A)\bigr]
  -
  \bigl[G(B\cup\{x\})-G(B)\bigr].
\]
If \(\mathrm{DR}(A,B,x)\ge 0\) for all such triples, then \(G\) is submodular.
\end{definition}

In the experiments, \(G_i(S)\) is evaluated on a paired seed \(i\) using common
random numbers: the base trajectory and scheduled trajectories share all
randomness except the corrector interventions.  This is a design choice for
variance reduction and causal comparability, not an additional theorem
assumption.

\section{Exact interaction decomposition}

The first mathematically clean contribution is to treat \(G\) as a finite-set
response surface.  No approximation is required.

\begin{theorem}[Möbius interaction decomposition]
Let \(G:2^{[T]}\to\R\) with \(G(\varnothing)=0\).  There exists a unique
collection of interaction coefficients \(\{I(U):\varnothing\ne U\subseteq[T]\}\)
such that
\[
  G(S)=\sum_{\varnothing\ne U\subseteq S} I(U)
  \qquad \text{for every } S\subseteq[T].
\]
They are given by Möbius inversion:
\[
  I(U)=\sum_{V\subseteq U} (-1)^{|U|-|V|}G(V).
\]
\end{theorem}

\begin{proof}
Define \(I(U)\) by the displayed inversion formula for nonempty \(U\), and set
\(I(\varnothing)=G(\varnothing)=0\) for this proof.  Then
\[
  \sum_{\varnothing\ne U\subseteq S} I(U)
  =
  \sum_{U\subseteq S}\sum_{V\subseteq U}(-1)^{|U|-|V|}G(V)
  =
  \sum_{V\subseteq S} G(V)
  \sum_{W\subseteq S\setminus V}(-1)^{|W|}.
\]
The inner sum equals \(1\) if \(V=S\) and \(0\) otherwise.  Hence the expression
equals \(G(S)\).  Uniqueness follows by induction on \(|S|\): once all
coefficients on strict subsets are known, \(I(S)=G(S)-\sum_{U\subsetneq S}I(U)\).
\end{proof}

The first-order coefficients are \(I(\{t\})=\Delta_t\).  The second-order
coefficients are \(I(\{s,t\})=\xi_{s,t}\).  Higher-order coefficients quantify
what remains unexplained by marginal and pair effects.

\subsection{Diagnostic ladder}

The decomposition gives a natural ladder:

\begin{center}
\begin{tabular}{lll}
\toprule
Dominant component & Statistical diagnosis & Policy implication \\
\midrule
First order \(I(\{t\})\) & marginal regime & separable ranking can work \\
Second order \(I(\{s,t\})\) & redundancy/complementarity & interaction diagnostics needed \\
Higher order \(I(U), |U|\ge3\) & non-low-order regime & true-\(G\) search only as probe \\
\bottomrule
\end{tabular}
\end{center}

This table is not a theorem about neural networks.  It is an exact
decomposition plus a diagnostic interpretation.

\section{Safe surrogate implications}

The next pieces are standard but useful.  They say when a diagnostic object
would justify a policy class.

\subsection{Marginal validity}

Define the additive surrogate
\[
  A(S):=\sum_{t\in S}\Delta_t.
\]

\begin{proposition}[Marginal regret under uniform additive residual]
Fix \(B\).  Suppose
\[
  \sup_{|S|=B}|G(S)-A(S)|\le \eta_B.
\]
Let \(S_A^\star\in\arg\max_{|S|=B}A(S)\) and
\(S_G^\star\in\arg\max_{|S|=B}G(S)\).  Then
\[
  G(S_G^\star)-G(S_A^\star)\le 2\eta_B.
\]
\end{proposition}

\begin{proof}
By the residual bound,
\[
G(S_G^\star)\le A(S_G^\star)+\eta_B
\le A(S_A^\star)+\eta_B
\le G(S_A^\star)+2\eta_B.
\]
\end{proof}

If \(\widehat{\Delta}_t\) is a proxy with
\(\sup_t|\widehat{\Delta}_t-\Delta_t|\le \varepsilon\), then selecting the
top-\(B\) times by \(\widehat{\Delta}_t\) adds at most \(2B\varepsilon\) regret
on \(A\), giving the familiar \(2\eta_B+2B\varepsilon\) bound.  This is the
right theorem for marginal rankers.  If the bound is vacuous or the empirical
rankers fail, the regime is not marginal.

\subsection{Pairwise validity}

Define the second-order surrogate
\[
  Q(S):=\sum_{t\in S}\Delta_t
  +\sum_{\{s,t\}\subseteq S}\xi_{s,t}.
\]

\begin{proposition}[Pairwise surrogate regret]
Suppose
\[
  \sup_{|S|=B}|G(S)-Q(S)|\le \zeta_B.
\]
Let \(S_Q^\star\in\arg\max_{|S|=B}Q(S)\).  Then
\[
  G(S_G^\star)-G(S_Q^\star)\le 2\zeta_B.
\]
If an estimated surrogate \(\widehat Q\) satisfies
\(\sup_{|S|=B}|Q(S)-\widehat Q(S)|\le\alpha_B\), and
\(\widehat S\) maximizes \(\widehat Q\), then
\[
  G(S_G^\star)-G(\widehat S)\le 2\zeta_B+2\alpha_B.
\]
\end{proposition}

\begin{proof}
The first statement is identical to the marginal proof with \(Q\) replacing
\(A\).  For the estimated case,
\[
Q(\widehat S)
\ge
\widehat Q(\widehat S)-\alpha_B
\ge
\widehat Q(S_Q^\star)-\alpha_B
\ge
Q(S_Q^\star)-2\alpha_B.
\]
Thus \(\widehat S\) has \(Q\)-optimization gap at most \(2\alpha_B\).  Apply
the exact-\(Q\) argument with this optimization gap.
\end{proof}

This is not a deep optimization theorem.  Its value is diagnostic: it states
precisely what must be true before a pairwise theory supports scheduling.

\section{Redundancy and diminishing returns}

The pairwise model gives a clean connection between negative interactions and
diminishing returns.

\begin{proposition}[Second-order submodularity criterion]
Let
\[
  Q(S)=\sum_{t\in S}a_t+\sum_{\{s,t\}\subseteq S}b_{s,t}.
\]
Then \(Q\) is submodular if and only if \(b_{s,t}\le 0\) for every pair
\(\{s,t\}\).
\end{proposition}

\begin{proof}
For \(A\subseteq B\) and \(x\notin B\),
\[
Q(A\cup\{x\})-Q(A)
=
a_x+\sum_{y\in A}b_{x,y},
\]
and similarly for \(B\).  Therefore
\[
\mathrm{DR}_Q(A,B,x)
=
-\sum_{y\in B\setminus A}b_{x,y}.
\]
This is nonnegative for all \(A,B,x\) if all \(b_{x,y}\le0\).  Conversely,
take \(A=\varnothing\), \(B=\{y\}\), \(x\ne y\).  Then
\(\mathrm{DR}_Q(\varnothing,\{y\},x)=-b_{x,y}\).  Submodularity implies
\(b_{x,y}\le0\).
\end{proof}

\begin{corollary}[Robustness to higher-order residuals]
Write \(G=Q+R\).  For any \(A\subseteq B\), \(x\notin B\),
\[
  \mathrm{DR}_G(A,B,x)
  =
  \mathrm{DR}_Q(A,B,x)
  +
  \mathrm{DR}_R(A,B,x).
\]
If \(\sup_S |R(S)|\le \rho\), then
\[
  |\mathrm{DR}_R(A,B,x)|\le 4\rho.
\]
Hence a positive second-order diminishing-returns gap larger than \(4\rho\)
survives the higher-order residual.
\end{corollary}

\begin{proof}
The first identity follows by linearity.  The second follows because
\(\mathrm{DR}_R\) is a signed sum of four \(R\)-values.
\end{proof}

This corollary is useful because it tells us why negative pair interactions
alone do not prove submodularity of \(G\).  They prove submodularity of the
second-order approximation.  To transfer the conclusion to \(G\), the
higher-order residual must be controlled.

\section{Corrector value as error, correctability, and downstream relevance}

The previous sections are exact finite-set/statistical statements.  They do
not explain why a corrector call is valuable.  For that, we introduce an
idealized Markov-chain layer.

At time \(t\), suppose there is an ideal local target \(\pi_t\) and current law
\(\nu_t\).  The corrector kernel \(K_t\) is useful if it contracts discrepancy
from \(\pi_t\).  Assume an entropy contraction condition:
\[
  \KL(\nu K_t\,\|\,\pi_t)
  \le
  (1-\kappa_t)\KL(\nu\,\|\,\pi_t),
  \qquad
  \kappa_t\in[0,1].
\]
This assumption is inspired by Markov-chain convergence theory, including
entropy contraction for Gibbs samplers under log-concavity
\cite{ascolani2024entropy}, but it is \emph{not} asserted for ProSeCo without
separate verification.

\begin{proposition}[Clean contraction-index statement]
If the local score of interest is negative relative entropy
\[
  \Phi_t(\nu):=-\KL(\nu\,\|\,\pi_t),
\]
and \(K_t\) satisfies the contraction condition above, then applying the
corrector improves this idealized score by at least
\[
  \Phi_t(\nu K_t)-\Phi_t(\nu)
  \ge
  \kappa_t\,\KL(\nu\,\|\,\pi_t).
\]
\end{proposition}

\begin{proof}
\[
\Phi_t(\nu K_t)-\Phi_t(\nu)
=
\KL(\nu\,\|\,\pi_t)-\KL(\nu K_t\,\|\,\pi_t)
\ge
\kappa_t \KL(\nu\,\|\,\pi_t).
\]
\end{proof}

This is the clean form of the intuition:
\[
  \text{ideal corrector value}
  \quad \gtrsim \quad
  \text{residual error}\times\text{local contraction strength}.
\]

For a terminal quality functional \(F\), one also needs downstream relevance.
Let \(h_t(x)\) be the expected terminal reward from state \(x\) at time \(t\)
under the remaining predictor/corrector plan.  If \(h_t\) is bounded, then
Pinsker's inequality gives
\[
  |(\nu K_t-\pi_t)h_t|
  \le
  \osc(h_t)\sqrt{\frac{1}{2}\KL(\nu K_t\,\|\,\pi_t)}
  \le
  \osc(h_t)\sqrt{\frac{1-\kappa_t}{2}\KL(\nu\,\|\,\pi_t)}.
\]
This does not prove the sign of the terminal improvement.  It says that a
contractive corrector reduces a bound on downstream bias relative to the ideal
local target, with strength governed by three factors:
\[
  \boxed{
  \text{residual error } \KL(\nu_t\|\pi_t)
  \quad
  \times
  \text{correctability } \kappa_t
  \quad
  \times
  \text{downstream sensitivity } \osc(h_t)
  }.
\]

This explains why uncertainty signals alone can fail.  High entropy can reflect
missing context rather than correctable error; a large residual is not useful
if \(K_t\) cannot contract it; and even a successful local correction has little
terminal value if downstream sensitivity is low.

\section{Connection to locally balanced proposals}

Zanella's locally balanced proposals are a theory for constructing informed
local MCMC moves in discrete spaces.  In a neighborhood \(\mathcal N(x)\) of
state \(x\), a locally informed proposal has the form
\[
  q_g(y\mid x)
  \propto
  g\!\left(\frac{\pi(y)}{\pi(x)}\right)K_0(x,y),
  \qquad y\in\mathcal N(x),
\]
where \(K_0\) is a base local proposal and \(g\) satisfies a balancing
condition of the form \(g(r)=r\,g(1/r)\).  Zanella shows that locally balanced
choices are optimal in a precise high-dimensional Peskun sense
\cite{zanella2020informed}.

The thesis should use this in the following disciplined way.

\begin{itemize}[leftmargin=1.5em]
\item Locally balanced proposals answer the \emph{kernel-design} question:
given a local target, how should the sampler propose local moves?
\item Informed correctors for discrete diffusion answer a related
generative-model question: how can a corrector use learned model information
to counter predictor approximation error?
\item Corrector timing asks a different question: given such a local kernel
\(K_t\), when should a finite number of calls to \(K_t\) be spent along an
inhomogeneous trajectory?
\end{itemize}

Thus the clean connection is:

\begin{quote}
An informed corrector is a locally informed Markov kernel.  Corrector timing is
the finite-budget scan problem of deploying such kernels along a generative
trajectory.
\end{quote}

We should not claim that ProSeCo is locally balanced unless this is separately
proved.  The theory only requires that \(K_t\) be a well-defined correction
kernel and that its statistical effect can be estimated.

\section{Connection to non-uniform scan Gibbs}

Random-scan Gibbs chooses which coordinate or block to update.  The selection
probabilities affect convergence.  Roberts and Sahu analyze how updating
schemes, blocking, and covariance structure affect Gibbs behavior
\cite{roberts1997updating}.  Chimisov, Latuszynski, and Roberts develop an
adaptive random-scan Gibbs sampler that adapts selection probabilities using a
spectral-gap criterion for a Gaussian analogue \cite{chimisov2018adapting}.

The corrector-timing analogy is:

\begin{center}
\begin{tabular}{ll}
\toprule
Gibbs / MCMC concept & Corrector-timing analogue \\
\midrule
coordinate or block & denoising time / correction opportunity \\
local update kernel & corrector kernel \(K_t\) \\
scan probability & timing policy or budget allocation \\
mixing improvement & terminal gain \(G(S)\) or local contraction index \\
adaptive scan & state-dependent corrector timing \\
\bottomrule
\end{tabular}
\end{center}

The analogy is useful but not literal.  A Gibbs sampler targets one stationary
distribution.  A masked-diffusion sampler is an inhomogeneous generative chain
whose intermediate targets change with time.  Therefore scan Gibbs supplies
vocabulary and motivation, not an off-the-shelf theorem.

In the additive idealization, if a single call at time \(t\) has value
\(c_t\), then allocating \(B\) deterministic calls to the largest \(c_t\)'s is
optimal.  The statistical content of this thesis is precisely that real
corrector calls need not remain additive: early calls change later states,
later values can become redundant, and pair/higher-order effects determine the
regime.

\section{Research question and theorem stack}

The recommended locked research question is:

\begin{quote}
For a fixed masked-diffusion predictor and fixed informed corrector, can
finite-budget corrector timing be statistically diagnosed from marginal
effects, conditional trajectory state, and interaction effects; and when is
timing reducible to marginal ranking versus interaction-driven correction?
\end{quote}

The corresponding theory stack is:

\begin{enumerate}[leftmargin=1.8em]
\item \textbf{Intervention setup.} Define \(J(S)\), \(G(S)\), \(\Delta_t\),
\(\tau_t(z)\), \(\xi_{s,t}\), and \(\mathrm{DR}(A,B,x)\).
\item \textbf{Exact interaction decomposition.} Use Möbius inversion to define
marginal, pairwise, and higher-order effects exactly.
\item \textbf{Safe low-order implications.} Prove marginal and pairwise
surrogate regret bounds under explicit finite-pool residual assumptions.
\item \textbf{Redundancy/submodularity diagnostic.} Show negative pair
coefficients imply submodularity of the second-order model, and quantify the
higher-order residual needed to transfer this conclusion to \(G\).
\item \textbf{Correctability index.} Under an idealized contraction assumption,
derive that local corrector value is governed by residual error, contraction
strength, and downstream sensitivity.
\item \textbf{Regime classification.} Classify a model/corrector/\(F\)/\(B\)
tuple as no-op, marginal, redundant/pairwise, locally searchable, or
higher-order/black-box.
\end{enumerate}

\section{How ProSeCo-OWT fits}

ProSeCo-OWT should be used as one empirical case study, not as universal
evidence.  The current diagnostics say:

\begin{itemize}[leftmargin=1.5em]
\item Corrector timing has positive headroom over uniform timing.
\item Tested separable rankers fail, so the regime is not purely marginal.
\item Pair interactions are structured and mostly negative, indicating
redundancy.
\item Cheap pairwise surrogates help at \(B=2\) but fail at \(B=3,4\), so
low-dimensional pairwise composition is insufficient.
\item Exact diminishing-return gaps are mildly positive on average, but not
strong enough for a clean submodularity claim.
\item One-swap neighborhoods contain many improvements, so local true-\(G\)
search is a useful probe.
\item Bayesian optimization remains unclear because the schedule-distance
kernel signal is weak.
\end{itemize}

Therefore ProSeCo-OWT is best described as an
\emph{interaction-driven, locally exploitable corrector-timing regime}.  This
is an empirical classification inside a model-agnostic framework.

\section{What is proved, what is borrowed, what is intuition}

\begin{center}
\begin{tabular}{p{0.30\linewidth}p{0.31\linewidth}p{0.28\linewidth}}
\toprule
Item & Status in this thesis & Source / inspiration \\
\midrule
Intervention estimands \(G,\Delta,\xi,\mathrm{DR}\) & Defined here; mathematically direct & potential-outcome/statistical response-surface viewpoint \\
Möbius decomposition & Exact theorem & finite Boolean-lattice inversion \\
Marginal and pairwise regret bounds & Exact under explicit sup residual assumptions & standard surrogate-regret argument \\
Second-order submodularity criterion & Exact for quadratic set functions & elementary set-function algebra \\
Contraction-index proposition & Exact under stated KL contraction assumption & Markov-chain contraction/Gibbs theory intuition \\
Locally balanced proposal connection & Conceptual bridge; not assumed for ProSeCo & Zanella informed proposals \\
Scan Gibbs connection & Conceptual bridge; not theorem transfer & Roberts--Sahu; adaptive scan Gibbs \\
ProSeCo-OWT classification & Empirical case study only & current diagnostics \\
\bottomrule
\end{tabular}
\end{center}

\section{Recommended thesis claim}

A careful final claim would be:

\begin{quote}
This thesis develops a statistical diagnostic framework for finite-budget
corrector timing in masked diffusion.  It treats corrector calls as
interventions in an inhomogeneous generative Markov chain, decomposes schedule
gain into marginal and interaction effects, and connects corrector timing to
the scan-design problem for locally informed Markov kernels.  On ProSeCo-OWT,
the framework identifies a non-marginal, redundant, locally exploitable regime:
separable marginal timing fails, pair interactions are mostly negative, weak
diminishing returns are present but not theorem-grade, and true-\(G\) search is
useful as an empirical probe.
\end{quote}

Claims to avoid:

\begin{itemize}[leftmargin=1.5em]
\item ProSeCo is locally balanced.
\item Corrector timing is submodular in general.
\item Bayesian optimization is justified by current diagnostics.
\item ProSeCo-OWT proves model-agnostic empirical behavior.
\item Search is the theoretical contribution.
\end{itemize}

\begin{thebibliography}{9}

\bibitem{zhao2024informed}
Yixiu Zhao, Jiaxin Shi, Feng Chen, Shaul Druckmann, Lester Mackey, and Scott
Linderman.
\newblock \emph{Informed Correctors for Discrete Diffusion Models}.
\newblock arXiv:2407.21243, 2024/2025.
\newblock \url{https://arxiv.org/abs/2407.21243}.

\bibitem{zanella2020informed}
Giacomo Zanella.
\newblock Informed proposals for local MCMC in discrete spaces.
\newblock \emph{Journal of the American Statistical Association},
115(530):852--865, 2020.
\newblock \url{https://arxiv.org/abs/1711.07424}.

\bibitem{roberts1997updating}
Gareth O. Roberts and Sujit K. Sahu.
\newblock Updating schemes, covariance structure, blocking and parameterization
for the Gibbs sampler.
\newblock \emph{Journal of the Royal Statistical Society: Series B},
59(2):291--317, 1997.
\newblock \url{https://ideas.repec.org/a/bla/jorssb/v59y1997i2p291-317.html}.

\bibitem{chimisov2018adapting}
Cyril Chimisov, Krzysztof Latuszynski, and Gareth Roberts.
\newblock \emph{Adapting the Gibbs Sampler}.
\newblock arXiv:1801.09299, 2018.
\newblock \url{https://arxiv.org/abs/1801.09299}.

\bibitem{ascolani2024entropy}
Filippo Ascolani, Hugo Lavenant, and Giacomo Zanella.
\newblock \emph{Entropy Contraction of the Gibbs Sampler under Log-Concavity}.
\newblock arXiv:2410.00858, 2024.
\newblock \url{https://arxiv.org/abs/2410.00858}.

\bibitem{secchi2025mwg}
Cecilia Secchi and Giacomo Zanella.
\newblock \emph{Spectral Gap of Metropolis-within-Gibbs under Log-Concavity}.
\newblock arXiv:2509.26175, 2025.
\newblock \url{https://arxiv.org/abs/2509.26175}.

\bibitem{lavenant2025}
Hugo Lavenant and Giacomo Zanella.
\newblock \emph{Error Bounds and Optimal Schedules for Masked Diffusions with
Factorized Approximations}.
\newblock arXiv:2510.25544, 2025.
\newblock \url{https://arxiv.org/abs/2510.25544}.

\end{thebibliography}

\end{document}
